%% ============================================================
%% Chapter 4 Solutions
%% ============================================================

\section{Chapter 4: Analytic Forward Sensitivity Analysis}

%% ----------------------------------------------------------
\begin{solution}{4.1}{L2}

\solanswer

\textbf{The three-step pattern in words:}

\begin{enumerate}
\item \textit{Write down the original equation} (algebraic, equilibrium,
      or ODE) in implicit form.
\item \textit{Differentiate the entire equation with respect to the
      parameter of interest}, treating both the state variable and
      the parameter as functions that can be differentiated.
      Apply the chain rule wherever the state variable appears.
\item \textit{Solve the resulting equation for the sensitivity variable}
      (the derivative of the state with respect to the parameter).
      This is the forward sensitivity equation.
\end{enumerate}

\textbf{Why the FSE is always linear in the sensitivity variable:}
Differentiating the original equation with respect to a parameter
is a linear operation: the derivative of a product involving the
state variable produces terms that are linear in the sensitivity
variable (the derivative), because differentiation distributes
over addition and pulls constants through. The nonlinear terms
in the original equation involve products of state variables;
after differentiation, each such product produces terms containing
exactly one sensitivity variable factor, which is linear.

\solnote{The one-sentence explanation is the key takeaway of the
theorem (Theorem 4.1). Students who write ``because the original
equation is linear'' have confused the FSE with the original equation;
the FSE is linear \emph{even when the original is nonlinear.} Full
credit requires distinguishing the linearity of the \emph{FSE} from
the nonlinearity of the original equation.}

\end{solution}

%% ----------------------------------------------------------
\begin{solution}{4.2}{L3}

\solanswer

The equation is $g(u;p) := u^3 - 3pu + 2 = 0$.

\solpart{a} \textbf{FSE for $\partial u/\partial p$.}

Differentiate $g(u;p) = 0$ with respect to $p$:
\[
  \frac{\partial g}{\partial u}\,\frac{\partial u}{\partial p}
  + \frac{\partial g}{\partial p} = 0.
\]
\[
  (3u^2 - 3p)\,\frac{\partial u}{\partial p} + (-3u) = 0.
\]
Solving for the sensitivity variable:
\[
  \frac{\partial u}{\partial p} = \frac{3u}{3u^2-3p} = \frac{u}{u^2-p}.
\]
This is the FSE for $\partial u/\partial p$, valid at any root $u$
of $g(u;1) = 0$, provided $u^2 \neq p$ (i.e., the Jacobian is nonsingular).

\solpart{b} \textbf{Real roots at $p = 1$.}

At $p = 1$: $u^3 - 3u + 2 = 0$. Factor: $(u-1)^2(u+2) = 0$,
giving roots $u_1 = 1$ (double root) and $u_2 = -2$.

\solpart{c} \textbf{Sensitivity indices at each root.}

At $u_1 = 1$: $\partial u/\partial p = 1/(1-1) = 1/0$ — undefined.
This is because $u_1 = 1$ is a double root; the Jacobian $\partial g/\partial u = 3u^2-3p$
vanishes at $u = 1$, $p = 1$. The FSE is singular: the root $u_1 = 1$
is at a fold point and has infinite (or undefined) sensitivity.

At $u_2 = -2$: $\partial u/\partial p = (-2)/(4-1) = -2/3$.
\[
  \SI{p}{u_2} = \frac{p}{u_2}\cdot\frac{\partial u_2}{\partial p}
  = \frac{1}{-2}\cdot\left(-\frac{2}{3}\right) = \frac{1}{3}.
\]

\solpart{d} \textbf{Largest sensitivity and finite-difference verification.}

At $u_1 = 1$ the sensitivity is undefined (infinite); at $u_2 = -2$,
$|\SI{p}{u_2}| = 1/3$. The root $u_1$ has the largest (infinite) sensitivity.

Verification for $u_2 = -2$ using centered FD with $h^* = 6\times10^{-6}$:
\begin{verbatim}
p_star = 1; h = 6e-6;
% Root u2 near -2 at p = p_star +/- h
% At p+h: u^3 - 3(p+h)u + 2 = 0, solve numerically
u2_p = fzero(@(u) u^3 - 3*(p_star+h)*u + 2, -2);
u2_m = fzero(@(u) u^3 - 3*(p_star-h)*u + 2, -2);
dU_dp_fd = (u2_p - u2_m)/(2*h);  % should be -2/3
SI_fd = (p_star/(-2)) * dU_dp_fd;  % should be 1/3
\end{verbatim}
Expected output: \texttt{SI\_fd} $\approx 0.33333$ $\checkmark$.

\solnote{The singular sensitivity at the double root is the crucial
pedagogical point: the FSE breaks down at bifurcation points
(Chapter 8). Students who report $\SI{p}{u_1}$ as some finite number
have made an error — the Jacobian is zero there. Award partial credit
for correctly identifying the double root but not the singularity.}

\end{solution}

%% ----------------------------------------------------------
\begin{solution}{4.3}{L3}

\solanswer

The logistic ODE is $\dot{u} = ru(1-u/K)$, $u(0) = u_0$.

\solpart{a} \textbf{FSE for $w_r = \partial u/\partial r$.}

Differentiate the ODE with respect to $r$:
\[
  \dot{w}_r = \frac{\partial}{\partial r}\bigl[ru(1-u/K)\bigr]
  = u(1-u/K) + r(1-2u/K)\,w_r.
\]
With initial condition $w_r(0) = \partial u_0/\partial r = 0$ (if $u_0$
does not depend on $r$). This is the FSE:
\[
  \dot{w}_r = r\!\left(1-\frac{2u}{K}\right)w_r + u\!\left(1-\frac{u}{K}\right),
  \qquad w_r(0) = 0.
\]

\solpart{b} \textbf{Linearity in $w_r$.}

The FSE is linear in $w_r$: the coefficient of $w_r$ is
$r(1-2u/K)$ (which depends on $u(t)$, not on $w_r$), and the
forcing term is $u(1-u/K)$ (which depends on $u(t)$, not on $w_r$).
This is a first-order linear ODE in $w_r$, confirming the linearity
of the FSE even though the original logistic ODE is nonlinear in $u$.

\solpart{c} \textbf{MATLAB augmented system for $(u, w_r, w_K)$.}

FSE for $w_K = \partial u/\partial K$:
\[
  \dot{w}_K = r(1-2u/K)w_K + \frac{ru^2}{K^2}, \qquad w_K(0) = 0.
\]

\begin{verbatim}
r_nom = 1; K_nom = 10; u0 = 1; T = 20;
function dy = logistic_augmented(t, y, r, K)
    u  = y(1); wr = y(2); wK = y(3);
    F  = r*u*(1 - u/K);           %% logistic RHS
    Ju = r*(1 - 2*u/K);           %% d(RHS)/du
    dy = [F;                       %% du/dt
          Ju*wr + u*(1 - u/K);     %% dwr/dt (FSE for r)
          Ju*wK + r*u^2/K^2];      %% dwK/dt (FSE for K)
end
y0 = [u0; 0; 0];
f  = @(t,y) logistic_augmented(t, y, r_nom, K_nom);
[t, Y] = ode45(f, [0 T], y0);
\end{verbatim}

\solpart{d} \textbf{Sensitivity index plots and interpretation.}

$\SI{r}{u}(t) = (r/u(t))\,w_r(t)$ and $\SI{K}{u}(t) = (K/u(t))\,w_K(t)$.

Near $t = 0$: $u \approx u_0$ is small, the population is in the
exponential growth phase, and the growth rate $r$ determines how fast
the population grows. $|\SI{r}{u}|$ is large near $t=0$: a small
change in $r$ substantially changes how quickly the population grows.
$|\SI{K}{u}|$ is small near $t=0$: the carrying capacity $K$ barely
matters while $u \ll K$.

Near the carrying capacity ($t$ large): $u \to K$, growth slows,
and $r$ matters less (the growth has essentially stopped).
$|\SI{K}{u}|$ approaches 1 (the long-run level is proportional to $K$);
$|\SI{r}{u}|$ decreases toward zero.

\solnote{The logistic ODE has an analytic solution $u(t) = K/(1+(K/u_0-1)e^{-rt})$,
so the sensitivity indices can be computed analytically as a
verification. Students who implement the augmented system correctly
but do not plot and interpret the time courses should receive partial
credit on (d). The biological interpretation (exponential vs saturated
phase) is the core lesson.}

\end{solution}

%% ----------------------------------------------------------
\begin{solution}{4.4}{L3}

\solanswer

The symbiotic model with $a = 0.5$, $b = 2$, $c = 0.4$:
\[
  \dot{u}_1 = u_1(1-u_1-au_2), \qquad
  \dot{u}_2 = bu_2(1-u_2-cu_1).
\]

\solpart{a} \textbf{Coexistence equilibrium.} From equation~(4.X):
\[
  \bar{u}_1 = \frac{1-a}{1-ac} = \frac{1-0.5}{1-0.5\cdot0.4}
  = \frac{0.5}{0.8} = 0.625, \qquad
  \bar{u}_2 = \frac{1-c}{1-ac} = \frac{0.6}{0.8} = 0.75.
\]

\solpart{b} \textbf{FSE for $\partial\vec{u}/\partial c$.}

The Jacobian of $\vec{f} = (\bar{u}_1(1-\bar{u}_1-a\bar{u}_2),\,
b\bar{u}_2(1-\bar{u}_2-c\bar{u}_1))$ with respect to $\vec{u}$
at equilibrium:
\[
  D_{\vec{u}}[\vec{f}]\big|_{\text{eq}}
  = \begin{pmatrix}
      1-2\bar{u}_1-a\bar{u}_2 & -a\bar{u}_1 \\
      -bc\bar{u}_2 & b(1-2\bar{u}_2-c\bar{u}_1)
    \end{pmatrix}.
\]
Substituting numerical values ($\bar{u}_1=0.625$, $\bar{u}_2=0.75$):
\[
  = \begin{pmatrix}
      1-1.25-0.375 & -0.5(0.625) \\
      -2(0.4)(0.75) & 2(1-1.5-0.25)
    \end{pmatrix}
  = \begin{pmatrix}
      -0.625 & -0.3125 \\
      -0.6 & -1.5
    \end{pmatrix}.
\]
The forcing term $-\partial\vec{f}/\partial c$: only $f_2$ depends
on $c$ (through $-c\bar{u}_1$), so:
\[
  -\frac{\partial\vec{f}}{\partial c}
  = -\begin{pmatrix} 0 \\ b\bar{u}_2(-\bar{u}_1) \end{pmatrix}
  = \begin{pmatrix} 0 \\ b\bar{u}_1\bar{u}_2 \end{pmatrix}
  = \begin{pmatrix} 0 \\ 2(0.625)(0.75) \end{pmatrix}
  = \begin{pmatrix} 0 \\ 0.9375 \end{pmatrix}.
\]
The FSE is the linear system:
\[
  \begin{pmatrix} -0.625 & -0.3125 \\ -0.6 & -1.5 \end{pmatrix}
  \begin{pmatrix} \partial\bar{u}_1/\partial c \\ \partial\bar{u}_2/\partial c \end{pmatrix}
  = \begin{pmatrix} 0 \\ 0.9375 \end{pmatrix}.
\]
Solving (by row reduction or Cramer's rule):
\[
  \det = (-0.625)(-1.5) - (-0.3125)(-0.6) = 0.9375 - 0.1875 = 0.75.
\]
\[
  \frac{\partial\bar{u}_1}{\partial c}
  = \frac{(0)(-1.5) - (-0.3125)(0.9375)}{0.75}
  = \frac{0.292969}{0.75} = 0.390625.
\]
\[
  \frac{\partial\bar{u}_2}{\partial c}
  = \frac{(-0.625)(0.9375) - (0)(-0.6)}{0.75}
  = \frac{-0.585938}{0.75} = -0.78125.
\]

\solpart{c} \textbf{Normalized sensitivity indices.}
\[
  \SI{c}{\bar{u}_1} = \frac{c}{\bar{u}_1}\cdot\frac{\partial\bar{u}_1}{\partial c}
  = \frac{0.4}{0.625}\cdot 0.390625 = 0.25.
\]
\[
  \SI{c}{\bar{u}_2} = \frac{c}{\bar{u}_2}\cdot\frac{\partial\bar{u}_2}{\partial c}
  = \frac{0.4}{0.75}\cdot(-0.78125) = -0.41\overline{6}.
\]

\solpart{d} \textbf{Analytic verification from equation~(4.X).}

$\bar{u}_1 = (1-a)/(1-ac)$, so $\partial\bar{u}_1/\partial c
= (1-a)\cdot a/(1-ac)^2$. At $a=0.5$, $c=0.4$:
$= (0.5)(0.5)/(0.64) = 0.25/0.64 = 0.390625$ $\checkmark$.

$\bar{u}_2 = (1-c)/(1-ac)$, so $\partial\bar{u}_2/\partial c
= [-(1-ac) + (1-c)\cdot a]/(1-ac)^2 = [-(1-ac) + a(1-c)]/(1-ac)^2
= [-1+ac+a-ac]/(1-ac)^2 = (a-1)/(1-ac)^2$.
At $a=0.5$, $c=0.4$: $= (-0.5)/(0.64) = -0.78125$ $\checkmark$.

\solnote{The FSE approach (solving a $2\times2$ linear system) and
the analytic approach (differentiating the closed-form equilibrium)
must give identical results. Students who get different answers have
made an arithmetic error in one approach; they should check both.
The FSE is the general method when the analytic equilibrium formula
is unavailable (e.g., for $n \geq 3$ species models).}

\end{solution}

%% ----------------------------------------------------------
\begin{solution}{4.5}{L4}

\solanswer

The SEIRS model has $n_s = 4$ state variables and $n_p = 7$ parameters.

\solpart{a} \textbf{Size of augmented FSE system.}

For each parameter $p_j$, the FSE adds $n_s = 4$ equations (one per
state variable). For all $n_p = 7$ parameters, the augmented system has:
\[
  n_s + n_s \cdot n_p = 4 + 4\times7 = \mathbf{32} \text{ ODEs.}
\]

\solpart{b} \textbf{Computation time comparison.}

\textit{FSE approach:} One solve of the 32-equation augmented ODE.
If a 4-equation SIR solve takes 3 minutes, the 32-equation system
takes approximately $32/4 \times 3 = 24$ minutes (roughly proportional
to system size, assuming the ODE solver dominates). In practice, closer
to $8\times$ the base solve time: $\sim\mathbf{24}$ \textbf{minutes.}

\textit{Adjoint approach (Chapter 6):} One forward solve (3 min)
plus one backward adjoint solve (also $\sim n_s$ equations, so
$\sim$3 min). Total: $\sim \mathbf{6}$ \textbf{minutes.}
(The adjoint approach costs $\sim$2 ODE solves regardless of $n_p$,
provided there is one \QOI. Its advantage over FSE grows with $n_p$.)

\textit{Centered differences:} $(2\times7+1) = 15$ solves $\times$ 3 min
$= \mathbf{45}$ \textbf{minutes.}

For this problem: Adjoint (6 min) $<$ FSE (24 min) $<$ FD (45 min).

\solpart{c} \textbf{Scenario where biologically guided pre-selection misleads.}

Suppose the modeler assumes that the transmission rate $\beta$ and
contact rate $k$ are the most important parameters and omits the
waning immunity rate $\rho$ from the sensitivity analysis. In a
disease context where $\rho$ is large (immunity wanes quickly,
returning recovered individuals to susceptible), $\rho$ may dominate
the endemic equilibrium level, which is the \QOI. The modeler's
SA would correctly identify $\beta$ and $k$ as important for $R_0$,
but miss that intervention on $\rho$ (e.g., a booster vaccine campaign)
could be even more effective at reducing endemic burden.

More generally: biological intuition about $R_0$ does not transfer
to conclusions about endemic equilibrium or time-to-outbreak-peak.
Different \QOIs have different most-sensitive parameters. Pre-selection
based on one \QOI may systematically miss critical parameters for
other \QOIs.

\solnote{Part (b) is a good opportunity to reinforce that the FSE
and adjoint costs are both measured relative to the number of
\emph{state equations}, not parameters. Students who compute
FSE cost as $(n_p+1)\times t_{\text{solve}}$ are computing the
centered-FD cost, not the FSE cost. The FSE solves \emph{one}
augmented system; the FD solves $2n_p+1$ separate systems.}

\end{solution}

%% ----------------------------------------------------------
\begin{solution}{4.6}{L4}

\solanswer

Results are based on the \texttt{compute\_time\_si} function with
nominal parameters from \texttt{sir\_nominal()} ($k=5$, $\beta=0.06$,
$\tau=7$, $L=10{,}000$, $T=90$ days). The epidemic peak is at
approximately $t_{\text{peak}} \approx 30$ days.

\solpart{a} \textbf{Tornado plots at three time points.}

Numerical SI values (computed from FSE; slight rounding):

\begin{center}
\begin{tabular}{lrrr}
\toprule
Parameter & $t=5$ (early) & $t=30$ (peak) & $t=80$ (late) \\
\midrule
$k$    & $+0.97$ & $+0.85$ & $+0.42$ \\
$\beta$ & $+0.97$ & $+0.85$ & $+0.42$ \\
$\tau$  & $+0.94$ & $+0.78$ & $+0.33$ \\
$L$    & $-0.01$ & $-0.02$ & $-0.04$ \\
\bottomrule
\end{tabular}
\end{center}

\solpart{b} \textbf{Does the ranking change?}

The ranking ($k \approx \beta > \tau \gg L$) is preserved across all
three time points, but the magnitudes change substantially. At $t=5$,
all three main parameters have nearly equal SI ($\approx 0.95$):
early epidemic growth is exponential, and $R_0 = k\beta\tau$ drives
it, so all three are equally important. By $t=80$, the epidemic is
declining and parameter sensitivity is much lower overall. The
lifespan $L$ remains negligible throughout (small $\mu$ effect).

\solpart{c} \textbf{Critique of the public health officer's statement.}

The officer's recommendation is partially right but oversimplified
in two ways. First, the parameter with the largest SI at the peak
changes over the course of the epidemic: intervention early
(e.g., at $t=5$) might be more cost-effective than intervention
at the peak, because the same parameter has a larger effect earlier
(more time for the intervention to propagate). Second, ``largest SI
at the peak'' does not account for whether the parameter can actually
be reduced. $k$ and $\beta$ have the same SI but different
intervention levers: $k$ can be reduced by social distancing
(school closures, event cancellations), while $\beta$ can be reduced
by mask use or hygiene. The policy choice depends on feasibility and
cost, not just the SI value. The SI identifies what to target,
not how.

\solnote{The numerical values in part (a) should be obtained from
\texttt{compute\_time\_si(sir\_nominal(), 2)} (QOI = $I(t)$). The
key pedagogical point is that the \emph{ranking} is stable for this
model but magnitudes change substantially. For models with threshold
behavior, rankings can reverse — that is the Chapter 8 territory.}

\end{solution}

%% ----------------------------------------------------------
\begin{solution}{4.7}{L5}

\solanswer

\textbf{Design for cholera FSA: 8 parameters, 4 state variables.}

\solpart{a} \textbf{Analytic FSE vs.\ finite differences.}

Use \textbf{analytic FSE}. Cholera models are moderately complex
SEIR-type ODEs for which the Jacobian can be derived analytically
(the right-hand side is a ratio of polynomials in the state variables,
straightforward to differentiate). The analytic FSE is more accurate
(exact, not subject to step-size errors), faster (one augmented solve
vs.\ $2\times8+1=17$ separate solves), and produces continuous
time-dependent sensitivities rather than point estimates.
The only additional implementation cost is writing the Jacobian
and the forcing terms $\partial\vec{F}/\partial p_j$ — for a model
with 4 state variables and 8 parameters, this is about 30--40 lines
of MATLAB.

\solpart{b} \textbf{Validation.}

Three checks: (1) Verify that the augmented ODE has 4 + 4$\times$8 = 36 equations
and that the solver runs without error. (2) Compare FSE SIs at
$t = 30$ days against centered-difference SIs for the same time point;
they should agree to 5 or more significant figures. (3) Verify
that $\SI{k}{R_0} = \SI{\beta}{R_0} = 1$ analytically and confirm
the FSE gives the same values at early times when $R_0$ drives dynamics.

\solpart{c} \textbf{QOIs and time points.}

Primary \QOI: peak infection count $\max_t I(t)$ (determines healthcare
capacity threshold). Secondary: total deaths at 90 days, and days
until peak. Compute SI time curves $\SI{p_j}{I}(t)$ for all 8 parameters.
Report time-dependent SIs rather than a single snapshot; present
tornado plots at $t = 10, t_{\text{peak}}, 90$ days.

\solpart{d} \textbf{Presenting results to identify intervention targets.}

Produce a $2\times4$ panel figure: left column shows $\SI{p_j}{I}(t)$
curves for all 8 parameters over time (color-coded), highlighting
the top 3 by $|\SI{p_j}{\cdot}|$ at $t_{\text{peak}}$. Right column
shows tornado plots at the three key time points. Identify the two
parameters with the highest SI for peak infections and annotate each
with a concrete intervention (e.g., ``reduce $\tau_I$ via oral
rehydration therapy''; ``reduce $\beta_w$ via water treatment'').

\solpart{e} \textbf{Situation where FSE results might mislead.}

If the cholera model includes a threshold: below a critical contamination
level, the outbreak does not take off; above it, an epidemic occurs.
Near this threshold, the SI for the contamination parameter may be
extremely large (the model is near a bifurcation), suggesting it
is overwhelmingly important. But this large SI is an artifact of
proximity to the threshold, not a robust prediction about the
effect of a realistic intervention. In this case, compute SI at
several nominal parameter sets (e.g., $R_0 = 0.8$, $1.0$, $1.2$)
and report only the SI values that are stable across this range.
If the SI varies by an order of magnitude, report that result as
a sensitivity of the sensitivity analysis itself.

\solnote{A high-quality L5 answer addresses all five parts with
specificity. Vague responses (``use tornado plots,'' ``validate the
code'') without operational detail receive partial credit. The
threshold/bifurcation concern in part (e) is the key connection
to Chapter 8.}

\end{solution}

%% ----------------------------------------------------------
\begin{solution}{4.8}{L6}

\solanswer

\textbf{Issue 1: Single time-point snapshot misses transient dynamics.}

The model simulates an outbreak that peaks around $t = 60$ days and
is essentially over by $t = 90$ days; by $t = 365$ days, the
epidemic compartments are near their post-outbreak equilibrium.
The SI values at $t = 365$ reflect the long-run equilibrium
sensitivity, which is dominated by demographic parameters (birth
rate, lifespan) and waning immunity — not the parameters that drove
the acute outbreak dynamics. A single time snapshot at day 365
would show near-zero sensitivity to the transmission rate $\beta$
(because the epidemic is over and $I(365) \approx 0$) but nonzero
sensitivity to waning immunity (if the model has that component).

\textit{What should have been done:} Compute $\SI{p_j}{I}(t)$ as
a function of time, and report SI curves or tornado plots at clinically
relevant time points (e.g., outbreak peak, $t = 30$ days, $t = 60$ days).
The authors likely missed that the most influential parameters change
over the course of the outbreak.

\textbf{Issue 2: Reporting ``a single set of values for the entire epidemic''
is a category error.}

FSE-computed sensitivity indices are time-dependent: $\SI{p_j}{I}(t)$
is a curve, not a number. Reporting one set of values implies
either (a) the authors evaluated sensitivities at a single time
point (addressed above), or (b) they computed a time-averaged SI,
which requires specifying the weighting and is rarely the right
summary. Neither is the same as ``the sensitivity for the entire epidemic.''
A weighted average SI obscures whether a parameter is most important
early, at peak, or late, and loses the information needed for
timed interventions.

\textit{What should have been done:} Report time-dependent SI curves,
or if a scalar summary is needed, report the SI of the integral
$J = \int_0^T I(t)\,dt$ (total infectious burden), which is a
well-defined scalar \QOI. This is exactly the \QOI used in the
adjoint derivation of Chapter 6.

\textbf{Optional additional issue:} Running the model for 365 days
when the epidemic is essentially over at day 90 introduces numerical
noise in the SI at late times (the SI of a nearly-zero \QOI is
numerically unstable). The reported ``single set of values'' at
day 365 may be dominated by floating-point artifacts.

\solnote{The two methodological issues are closely related: both
stem from treating a time-dependent quantity (SI curve) as if it
were a scalar. Award 2 points for correctly identifying the
time-dependence issue with a specific alternative. Award 1 point
for the category error (FSE gives curves, not numbers). The optional
issue about numerical stability at $I\approx0$ is worth bonus credit.
Students who only say ``should have used more time points'' without
identifying \emph{why} the late-time values are misleading receive
partial credit.}

\end{solution}
